\documentclass[10pt,a4paper]{amsart}
\usepackage[margin=19mm]{geometry}
\usepackage{amsmath,amssymb,mathtools}
\usepackage[scr=rsfs,cal=euler]{mathalfa}
\usepackage{xfrac}
\usepackage[unicode,hidelinks,bookmarks=false]{hyperref}
\newcommand{\ie}{i.e.\ }
\newcommand{\cf}{cf.\ }
\newcommand{\resp}{respectively\ }
\newcommand{\Subsection}{\subsection}
\providecommand{\nobreakdash}{\nobreak\hbox{-}}
\setlength{\emergencystretch}{2em}
\multlinegap0pt
\usepackage{bm}
\usepackage{bbm}
\usepackage{dsfont}
\usepackage{xspace}
\usepackage{cancel}
\usepackage{mathtools}
\usepackage{amsthm}
\makeatletter
\@ifundefined{theorem}{\newtheorem{theorem}{Theorem}}{}
\@ifundefined{prop}{\newtheorem{prop}[theorem]{Proposition}}{}
\makeatother
\title[Sign regularity preserving linear operators]{Sign regularity preserving linear operators}
\author{Projesh Nath Choudhury, Shivangi Yadav}
\date{Source: arXiv:2408.02428v3 (CC BY 4.0)}
\begin{document}
\maketitle

\noindent\emph{Source and licence.} Sign regularity preserving linear operators. Version arXiv:2408.02428v3 (CC BY 4.0), distributed under the Creative Commons Attribution 4.0 International licence, as recorded in the arXiv licence metadata for this exact version and shown on the abstract page https://arxiv.org/abs/2408.02428v3. Full rights notice: licenses/arxiv-2408.02428-CC-BY-4.0.txt.\par\smallskip

\noindent\textit{Editorial scope.} Counted unit: the Prop whose source label is \texttt{SR\_rectangle\_S11Smn\_S1nSm1\_combo} in the article (source0.tex lines 537--546, source label \texttt{SR\_rectangle\_S11Smn\_S1nSm1\_combo}), delivered together with the article's own complete original proof of it (source0.tex lines 548--591, one \texttt{proof} environment of 44 lines). No mathematics is written, repaired or re-derived: statement and proof are the article's own text line for line. Required context retained uncounted: basis (lines 226--228); SR\_square\_monomial (lines 273--275); SR\_square\_S11Snn\_S1nSn1\_combo (lines 326--335); SR\_rectangle\_S11SmnS1nSm1\_elements (lines 485--487); Proposition\_L(J)\_rank\_1 (lines 513--515). Every definition, display and label of those lines is retained unchanged. Omitted from this package and not redistributed: 1--225, 229--272, 276--325, 336--484, 488--512, 516--536, 547--547, 592--1400. Nothing inside a retained excerpt is omitted or altered: every retained body is delivered exactly as the article's lines are written, so no omission receipt of any kind accompanies this package. Citations: the retained excerpts carry no citation command at all, so no bibliography entry of the article is transcribed and no reference is redistributed. Reference adaptation: none was needed and none was made; every reference inside the delivered unit resolves inside it, so no unresolved reference or citation is produced. Printed slips kept literal: no printed word, symbol, hypothesis or number of the counted statement or of its proof was altered. Layout and class: the article's own class is \texttt{amsart} with packages including geometry, eucal, hyperref, xcolor, enumerate, amsmath, amsthm, amsfonts, amssymb, tocvsec2, mathdots, mathtools, caption, subcaption; the wrapper is the lane's pinned \texttt{amsart} editorial wrapper at 10pt on A4, which restates the theorem-like environments the retained text needs and adds the article's own macro definitions that the retained text needs (none), with extra wrapper packages (bm, bbm, dsfont, xspace, cancel, mathtools). The delivered numbering is the unit's own. Assets: no retained span contains a figure environment, a graphics-inclusion command, a PDF-page inclusion, a table or a TikZ picture; no image file is copied into this package, the package declares no asset dependency and no image file is redistributed with it. Licence: version arXiv:2408.02428v3 of this article is distributed under the Creative Commons Attribution 4.0 International licence, as recorded in the arXiv licence metadata for this exact version and shown on the abstract page https://arxiv.org/abs/2408.02428v3. A full rights notice is delivered at licenses/arxiv-2408.02428-CC-BY-4.0.txt. Characters: every retained excerpt is pure ASCII, so the delivered engine, XeLaTeX with its default Latin Modern text font, reports no missing character.
		\begin{equation}\label{basis}
		\mathfrak{B}=(E_{11},\ldots,E_{kk};E_{12},\ldots,E_{mn})
		\end{equation} a fixed ordered basis of each $\mathbb{R}^{m\times n}$, where $k=\min\{m,n\}$. The basis elements are arranged such that $E_{ii}$ appear first in increasing order of $i$, followed by the remaining elements arranged in lexicographic order.

\begin{prop}\label{SR_square_monomial}
	Let $L\in\mathbb{R}^{mn\times mn}$ be the matrix representation of $\mathcal{L}$ with respect to the ordered basis $\mathfrak{B}$ as in~\eqref{basis}. Then $L$ is a monomial matrix.
\end{prop}

	\begin{prop}\label{SR_square_S11Snn_S1nSn1_combo}
		For $\mathcal{L}\in P(\mathcal{SR}_2)$, the following pairwise combinations are possible.
		\begin{itemize}
			\item[(i)] $S_{11}=\{E_{11}\}\ \text{and} \ S_{nn}=\{E_{nn}\}, \ \text{or} \ S_{11}=\{E_{nn}\} \ \text{and}\ S_{nn}=\{E_{11}\}, \ \text{or}$ \\
			$S_{11}=\{E_{1n}\}\ \text{and}\ S_{nn}=\{E_{n1}\}, \ \text{or}\ S_{11}=\{E_{n1}\} \ \text{and}\ S_{nn}=\{E_{1n}\}$.
			
			\item[(ii)] $S_{1n}=\{E_{1n}\} \ \text{and}\ S_{n1}=\{E_{n1}\}, \ \text{or}\ S_{1n}=\{E_{n1}\} \ \text{and}\ S_{n1}=\{E_{1n}\}, \ \text{or}$ \\
			$S_{1n}=\{E_{11}\} \ \text{and}\ S_{n1}=\{E_{nn}\}, \ \text{or}\ S_{1n}=\{E_{nn}\} \ \text{and}\ S_{n1}=\{E_{11}\}$.
		\end{itemize}
	\end{prop}

 \begin{prop}\label{SR_rectangle_S11SmnS1nSm1_elements}
		For $\mathcal{L}\in P(\mathcal{SR}_2)$, the element in each of the sets $S_{11}$, $S_{mn}$, $S_{1n}$, and $S_{m1}$ must be one of $E_{11}$, $E_{mn}$, $E_{1n}$, $E_{m1}$.
	\end{prop}

	\begin{prop}\label{Proposition_L(J)_rank_1}
		Let $m>n\geq2$. For $\mathcal{L}\in P(\mathcal{SR}_2)$, the matrix $\mathcal{L}(J_{m\times n})$ has rank one.
	\end{prop}

	\begin{prop}\label{SR_rectangle_S11Smn_S1nSm1_combo} 
		For $\mathcal{L}\in P(\mathcal{SR}_2)$, the following pairwise combinations are possible.
		\begin{itemize}
			\item[(i)] $S_{11}=\{E_{11}\} ~\text{and}~ S_{mn}=\{E_{mn}\}, ~\text{or}~ S_{11}=\{E_{mn}\} ~\text{and}~ S_{mn}=\{E_{11}\}, ~\text{or}$ \\
			$S_{11}=\{E_{1n}\} ~\text{and}~ S_{mn}=\{E_{m1}\}, ~\text{or}~ S_{11}=\{E_{m1}\} ~\text{and}~ S_{mn}=\{E_{1n}\}$.
			
			\item[(ii)] $S_{1n}=\{E_{1n}\} ~\text{and}~ S_{m1}=\{E_{m1}\},~\text{or}~ S_{1n}=\{E_{m1}\} ~\text{and}~ S_{m1}=\{E_{1n}\}, ~\text{or}$ \\
			$S_{1n}=\{E_{11}\} ~\text{and}~ S_{m1}=\{E_{mn}\},~\text{or}~ S_{1n}=\{E_{mn}\} ~\text{and}~ S_{m1}=\{E_{11}\}$.
		\end{itemize}
	\end{prop}

	\begin{proof} First, we prove (i). The proof is similar to that
	of Proposition~\ref{SR_square_S11Snn_S1nSn1_combo}, except the $n=2$ case. From Proposition~\ref{SR_rectangle_S11SmnS1nSm1_elements}, we have that $S_{11}$ equals one of the sets $\{E_{11}\}, \{E_{mn}\}, \{E_{m1}\}, \{E_{1n}\}$ and similarly, $S_{mn}$ is one of $\{E_{11}\}, \{E_{mn}\}, \{E_{m1}\}, \{E_{1n}\}$. Since $L$ is a  monomial matrix, out of sixteen possible combinations of $S_{11}$ and $S_{mn}$, the four cases wherein $S_{11}=S_{mn}$ are discarded straightaway. We next show that the following eight other combinations of $S_{11}$ and $S_{mn}$ are also not possible:
	\begin{align*}
		&S_{11}=\{E_{m1}\} ~\text{and}~ S_{mn}=\{E_{mn}\}, \quad 
		S_{11}=\{E_{11}\} ~\text{and}~ S_{mn}=\{E_{1n}\}, \\
		&S_{11}=\{E_{mn}\} ~\text{and}~ S_{mn}=\{E_{m1}\}, \quad 
		S_{11}=\{E_{mn}\} ~\text{and}~ S_{mn}=\{E_{1n}\}, \\
		&S_{11}=\{E_{m1}\} ~\text{and}~ S_{mn}=\{E_{11}\}, \quad
		S_{11}=\{E_{11}\} ~\text{and}~ S_{mn}=\{E_{m1}\}, \\
		&S_{11}=\{E_{1n}\}~\text{and}~ S_{mn}=\{E_{11}\}, \quad
		S_{11}=\{E_{1n}\}~\text{and}~ S_{mn}=\{E_{mn}\}.
	\end{align*}
	Suppose $S_{11}=\{E_{m1}\}$ and $S_{mn}=\{E_{mn}\}$. Let $n=2$ and for $c>0$, define
		\[J(c):= \begin{pmatrix}
			1+c \qquad c \\
			J_{(m-1) \times 2}
		\end{pmatrix}\in \mathcal{SR}\implies \mathcal{L}(J(c))\in \mathcal{SR}.\]
	By Propositions~\ref{SR_square_monomial}~and~\ref{SR_rectangle_S11SmnS1nSm1_elements}, $\mathcal{L}(J(c))$ can be either of the following:
	\begin{align*}
		\begin{pmatrix}
			y_{11} & cy_{12} \\
			y_{21} & y_{22} \\
			\vdots & \vdots \\
			(1+c)y_{m1} & y_{m2}
		\end{pmatrix}~\text{or}~\begin{pmatrix}
			cy_{11} & y_{12} \\
			y_{21} & y_{22} \\
			\vdots & \vdots \\
			(1+c)y_{m1} & y_{m2}
		\end{pmatrix}.
	\end{align*}
	For the first case, using the fact that $\mathcal{L}(J_{m\times2})$ has rank 1 by Proposition~\ref{Proposition_L(J)_rank_1}, the following two $2\times2$ minors of $Y(c)$ have opposite signs for an appropriate choice of $c$:
		\[\det\begin{pmatrix}
			y_{11} & cy_{12} \\
			y_{21} & y_{22} 
		\end{pmatrix}>0~\text{and}~
		\det\begin{pmatrix}
			y_{21} & y_{22} \\
			(1+c)y_{m1} & y_{m2} 
		\end{pmatrix}<0.\]
	For the other case, we can choose $c$ large enough such that we obtain two minors of size $2\times2$ of $Y(c)$ having opposite signs. Hence, we conclude that $S_{11}\neq\{E_{m1}\}$ and $S_{mn}\neq\{E_{mn}\}$ for $n=2$. Now the proof for $m>n\geq3$ follows similarly to the proof of Proposition~\ref{SR_square_S11Snn_S1nSn1_combo}.
	
	Using a similar argument as in the above proof shows that (ii) holds.
	\end{proof}

\end{document}
